\subsection{\titlemath{$\UT{1}$}{UT₁}: Scaling Complexity Intuition}

\aside{
    Although we will rigorously investigate UT's scaling complexity in \autoref{sec:ut-complexity}, it will prove useful to introduce $N_1$ here and to gain an intuition for UT's scalability before \autoref{sec:practical-considerations}.
}

How scalable is \UT1?

Let us start by assuming that a single blockchain has $O(c)$ capacity (measured in transactions per second) and is thus $O(c)$ scalable.
We'll also assume that the \emph{headers} are $O(1)$ bytes (i.e., a constant size).

Note that a single blockchain is a 0-simplex.

A 1-simplex has two blockchains, and each chain must record the headers of the other and the respective PoRs.
For the sake of simplicity, let's assume that the PoRs are also $O(1)$.
Thus, each chain loses $O(1)$ of its $O(c)$ capacity, and the network's total capacity is $2O(c)-2O(1) = O(c)$.

Since each chain has finite capacity, there is an absolute maximum size to a simplex.
We can keep adding more chains until all blocks are 100\% full of headers and PoRs.
However, this system has 0 capacity for transactions.
So, there must be a capacity maxima with more than 2 simplex-chains and fewer than the absolute maximum.

What we're observing is that the overall capacity is the \emph{product} of the size of the simplex (in chains) and each chain's remaining capacity --- a quadratic relationship.

Let's call the proportion of capacity for transactions $t$, and the proportion for reflections $r$.
We know that $t + r = 1$, since $t+r$ is 100\% of a chain's capacity.
However, the \emph{network-wide} capacity is proportional to the \emph{product} of $t$ and $r$.
Maximizing $t r$ suggests that a chain's capacity should be split approximately 50/50 between $t$ and $r$, so $t=r=\half=O(1)$.
Given that $t$ and $r$ are proportions of the capacity, we can say $O(ct) + O(cr) = O(c)$.
Thus, the overall capacity is given by $O(ct) \cdot O(cr) = O(c^2)$.

This relationship is analogous to the relationship between a rectangle's maximum area for a given perimeter length: the area scales with the square of the perimeter (and such a rectangle is a square).

Let's say $N_1$ is the number of chains which maximizes a simplex's capacity (for transactions).
To get a feel for $N_1$, let's estimate it using example chains based on modified Bitcoin parameters.

First, we'll add a merkle root of headers and PoRs to the header, increasing header size from 80 to 112 bytes.
Second, we'll reduce the block frequency down from 10 minutes to 1 minute, and scale the block limit down from 1 MB to 100 kB (this will give us more realistic numbers).
Let's assume the merkle branch leading to a header is 256 bytes, so the total PoR overhead is 368 B / block / chain.
We want to use 50\% of the overall capacity (100 kB / block) for PoRs, so each block should contain \sim 50 kB worth of PoRs.
Dividing the PoR capacity by the per-chain overhead gives us $N_1 \approx 139$ chains and a network-wide capacity equivalent to \sim 70 Bitcoin networks.
